%%%%%%%%%%%%%%%%%
% This is an sample CV template created using altacv.cls
% (v1.7.4, 30 July 2025) written by LianTze Lim (liantze@gmail.com). Compiles with pdfLaTeX, XeLaTeX and LuaLaTeX.
%
%% It may be distributed and/or modified under the
%% conditions of the LaTeX Project Public License, either version 1.3
%% of this license or (at your option) any later version.
%% The latest version of this license is in
%%    http://www.latex-project.org/lppl.txt
%% and version 1.3 or later is part of all distributions of LaTeX
%% version 2003/12/01 or later.
%%%%%%%%%%%%%%%%

%% Use the "normalphoto" option if you want a normal photo instead of cropped to a circle
% \documentclass[10pt,a4paper,withhyper,normalphoto]{altacv}

\documentclass[10pt,a4paper,withhyper]{altacv}
%% AltaCV uses the fontawesome5 and simpleicons packages.
%% See http://texdoc.net/pkg/fontawesome5 and http://texdoc.net/pkg/simpleicons for full list of symbols.

% Change the page layout if you need to
\geometry{left=1.25cm,right=1.25cm,top=1.5cm,bottom=1.5cm,columnsep=1.2cm}

% The paracol package lets you typeset columns of text in parallel
\usepackage{paracol}

% Change the font if you want to, depending on whether
% you're using pdflatex or xelatex/lualatex
% WHEN COMPILING WITH XELATEX PLEASE USE
% xelatex -shell-escape -output-driver="xdvipdfmx -z 0" sample.tex
\iftutex
  % If using xelatex or lualatex:
  \setmainfont{Roboto}
  \setsansfont{roboto}
  \renewcommand{\familydefault}{\sfdefault}
\else
  % If using pdflatex:
  \usepackage[rm,defaultsans]{roboto}
  \usepackage[defaultsans]{roboto}
  % \usepackage{sourcesanspro}
  \renewcommand{\familydefault}{\sfdefault}
\fi
\usepackage[sfdefault]{roboto} % Sets Roboto as the sans-serif font
\usepackage[T1]{fontenc}       % Helps LaTeX find bold/italic variants
% Change the colours if you want to
\definecolor{SlateGrey}{HTML}{2E2E2E}
\definecolor{LightGrey}{HTML}{666666}
\definecolor{DarkPastelRed}{HTML}{450808}
\definecolor{PastelRed}{HTML}{8F0D0D}
\definecolor{GoldenEarth}{HTML}{E7D192}
\colorlet{name}{black}
\colorlet{tagline}{PastelRed}
\colorlet{heading}{black}
\colorlet{headingrule}{PastelRed}
\colorlet{subheading}{PastelRed}
\colorlet{accent}{black}
\colorlet{emphasis}{SlateGrey}
\colorlet{body}{LightGrey}

% Change some fonts, if necessary
\renewcommand{\namefont}{\Huge\rmfamily\roboto}
\renewcommand{\personalinfofont}{\footnotesize}
\renewcommand{\cvsectionfont}{\LARGE\rmfamily\roboto}
\renewcommand{\cvsubsectionfont}{\large\bfseries}


% Change the bullets for itemize and rating marker
% for \cvskill if you want to
\renewcommand{\cvItemMarker}{{\small\textbullet}}
\renewcommand{\cvRatingMarker}{\faCircle}
% ...and the markers for the date/location for \cvevent
% \renewcommand{\cvDateMarker}{\faCalendar*[regular]}
% \renewcommand{\cvLocationMarker}{\faMapMarker*}


% If your CV/résumé is in a language other than English,
% then you probably want to change these so that when you
% copy-paste from the PDF or run pdftotext, the location
% and date marker icons for \cvevent will paste as correct
% translations. For example Spanish:
% \renewcommand{\locationname}{Ubicación}
% \renewcommand{\datename}{Fecha}


%% Use (and optionally edit if necessary) this .tex if you
%% want to use an author-year reference style like APA(6)
%% for your publication list
% \input{pubs-authoryear}

%% Use (and optionally edit if necessary) this .tex if you
%% want an originally numerical reference style like IEEE
%% for your publication list
\input{pubs-num}

%% sample.bib contains your publications
\addbibresource{sample.bib}








\begin{document}
\name{\textbf{Jahid Howlader}}
\tagline{Student - Aspiring Software Engineer}
%% You can add multiple photos on the left or right
\photoR{3.0cm}{image/photo.jpg}
% \photoL{2.5cm}{Yacht_High,Suitcase_High}

\personalinfo{%
  % Not all of these are required!
  \email{jahidhowlader2006@gmail.com}
  \phone{+49 162 1651752}

  \location{Berlin, Germany}
  \homepage{Jahid Howlader}\href{https://jahidhowlader6.github.io/}{}
  % \twitter{@twitterhandle}
  \linkedin {Jahidhowlader6}\href{https://www.linkedin.com/in/jahidhowlader6/?skipRedirect=true}{}  
  \github{Jahidhowlader6}\href{https://github.com/jahidhowlader6}{}
  
  %% You can add your own arbitrary detail with
  %% \printinfo{symbol}{detail}[optional hyperlink prefix]
  % \printinfo{\faPaw}{Hey ho!}[https://example.com/]

  %% Or you can declare your own field with
  %% \NewInfoFiled{fieldname}{symbol}[optional hyperlink prefix] and use it:
  
  %%
  %% For services and platforms like Mastodon where there isn't a
  %% straightforward relation between the user ID/nickname and the hyperlink,
  %% you can use \printinfo directly e.g.
  % \printinfo{\faMastodon}{@username@instace}[https://instance.url/@username]
  %% But if you absolutely want to create new dedicated info fields for
  %% such platforms, then use \NewInfoField* with a star:
  \NewInfoField*{mastodon}{\faMastodon}
  %% then you can use \mastodon, with TWO arguments where the 2nd argument is
  %% the full hyperlink.
  
}

\makecvheader
%% Depending on your tastes, you may want to make fonts of itemize environments slightly smaller
% \AtBeginEnvironment{itemize}{\small}

%% Set the left/right column width ratio to 6:4.
\columnratio{0.6}

% Start a 2-column paracol. Both the left and right columns will automatically
% break across pages if things get too long.
\begin{paracol}{2}
\cvsection{\textbf{Experience}}

\cvevent{Barista}{Creepy Bh}{1 month}{Riffa, Bahrain}
\begin{itemize}
\item Made cookies and ice cream
\item Made different types of coffees
\end{itemize}

\divider

\cvevent{Freelance designer}{RedBubble}{2022-2024}{Bahrain}
\begin{itemize}
\item Created designed and managed and online shop
\item Generated 150+ sales
\end{itemize}

\cvsection{\textbf{Projects}}

\cvevent{Avion}{Passion Project}{Ongoing}{}
\begin{itemize}
A web-app that lets aviation enthusiasts share their passion for aviation. This app will provide a curated feed of aviation content the user is interested in.
\end{itemize}
\href{https://docs.google.com/document/d/1pWOtEOxm7aP5mE0FJZz1KgR_11sMZT2RKtRK7JkRwXE/edit?usp=sharing}{\textbf{\textit{\textcolor{blue}{Click here to see}}}}
\divider

\cvevent{ASU Hackathon}{ASU, Bahrain}{2 days}{}
\begin{itemize}A Cybersecurity competition with University and other school students, hosted by \href{https://www.asu.edu.bh/}{\textbf{\textit{Applied Science University, Bahrain}}} and \href{https://engage.isaca.org/bahrainchapter/home}{\textit{\textbf{ISACA Bahrain Chapter}}}.
\end{itemize}

\href{https://www.instagram.com/p/DIjhPNSsdv_/?utm_source=ig_web_copy_link&igsh=MzRlODBiNWFlZA==}{\textbf{\textit{\textcolor{blue}{Click here to see}}}}
\medskip



% use ONLY \newpage if you want to force a page break for
% ONLY the current column


\switchcolumn

\cvsection{\textbf{About Me}}

\begin{quote}
I'm an aspiring software engineer with a passion for Creativity and Aviation. I love building things!
\end{quote}

\cvsection{\textbf{Skills/Hobbies}}

\cvachievement{\faTrophy}{Creative}{Photography, Video editing, Photo editing, Graphic design, sketching}

\divider

\cvachievement{\faWifi}{Tech}{Knowledge of Computer Hardware, Software, AI tools, and New tech}

\divider

\cvachievement{\faIcons }{Other}{watching documentaries and investigative journalism, and interested in Si-Fi movies, shows and games  }

\cvsection{\textbf{Strengths}}

% Don't overuse these \cvtag boxes — they're just eye-candies and not essential. If something doesn't fit on a single line, it probably works better as part of an itemized list (probably inlined itemized list), or just as a comma-separated list of strengths.

\cvtag{Hard-working}
\cvtag{Eye for detail}\\
\cvtag{Motivator \& Leader}
\cvtag{Creativity}

\divider\smallskip

\cvtag{Front-End development}
\cvtag{Photography}\\
\cvtag{Video Editing}
\cvtag{Engineering solutions} \\
\cvsection{\textbf{Languages}}

\cvskill{English}{5}
\divider

\cvskill{Hindi, Urdu, Bangla}{4}
\divider

\cvskill{German}{2} %% Supports X.5 values.

%% Yeah I didn't spend too much time making all the
%% spacing consistent... sorry. Use \smallskip, \medskip,
%% \bigskip, \vspace etc to make adjustments.
\medskip
\cvsection{\textbf{Programming Languages}}

\cvskill{HTML/CSS}{4}
\divider
\cvskill{AI tools}{4}
\divider
\cvskill{LaTeX}{3}
\divider
\cvskill{C++}{2} %% Supports X.5 values.
\divider

\cvskill{Python}{1}

%% Yeah I didn't spend too much time making all the
%% spacing consistent... sorry. Use \smallskip, \medskip,
%% \bigskip, \vspace etc to make adjustments.
\medskip
\switchcolumn

\cvsection{\textbf{Education}}

\cvevent{Fbise HSSC-II}{Pakistan Urdu School, Bahrain}{Sept 2002 -- June 2006}{}
Calculus, Trigonometry, Physics, English, C++, MS Excel, MS Access, HTML, Computer Science (Network architecture, Hardware etc). \\
\href{https://drive.google.com/file/d/1UKWIip42YEXn_VlNceEmjSV-ZG-pNDXX/view?usp=sharing}{\textbf{\textit{\textcolor{blue}{Click here to See}}}}

\divider

\cvevent{BSc. Software Engineering}{GISMA University of Applied Science, Berlin}{April 2026-Ongoing}{}
\href{https://drive.google.com/file/d/1kEWC-4LXxUv5OTvcrJ9Dxp2-BcYU5JIO/view?usp=sharing}{\textit{\textcolor{blue}{\textbf{Click here to see}}}}    
\divider


\end{paracol}


\end{document}
